\documentclass{article}

\usepackage{booktabs}
\usepackage{tabularx}
\usepackage{hyperref}
\usepackage{graphicx}

\hypersetup{
    colorlinks=true,
    linkcolor=red,
    citecolor=green,
    filecolor=magenta,
    urlcolor=cyan
}

\title{Software Requirements Specification\\\progname}

\author{\authname}

\date{}

\input{../Comments.text}
\input{../Common.text}

\begin{document}

\maketitle
\thispagestyle{empty}

~\newpage

\pagenumbering{roman}

\begin{table}[hp]
\caption{Revision History} \label{TblRevisionHistory}
\begin{tabularx}{\textwidth}{llX}
\toprule
\textbf{Date} & \textbf{Developer(s)} & \textbf{Change}\\
\midrule
Date1 & Name(s) & Description of changes\\
Date2 & Name(s) & Description of changes\\
... & ... & ...\\
\bottomrule
\end{tabularx}
\end{table}

~\newpage

\tableofcontents

~\newpage

\pagenumbering{arabic}

% Structure follows Bertrand Meyer's Handbook of Requirements and Business
% Analysis (four books: Goals, Environment, System, Project), adapted to LaTeX
% for this course deliverable.

\section{Goals} \label{sec:goals}

The Goals book addresses stakeholders broadly: why the project exists, what
benefits are expected, and what is in or out of scope.

\subsection{Context and Overall Objectives} \label{sec:g1}

\wss{High-level view of the project: organizational context and reason for
building a system.}

\subsection{Current Situation} \label{sec:g2}

\wss{Current state of the processes the project and resulting system will
address.}

\subsection{Expected Benefits} \label{sec:g3}

\wss{New or improved processes made possible by the project's results. Core of
the Goals book.}

\subsection{Functionality Overview} \label{sec:g4}

\wss{Short overview of principal system behaviours; details belong in
Section~\ref{sec:system}.}

\subsection{High-Level Usage Scenarios} \label{sec:g5}

\wss{Main usage patterns in user terms only (no special/error cases). Detailed
scenarios belong in Section~\ref{sec:s4}.}

\subsection{Limitations and Exclusions} \label{sec:g6}

\wss{What the system will not address. Risks belong in Section~\ref{sec:p6}, not
here.}

\subsection{Stakeholders and Requirements Sources} \label{sec:g7}

\wss{Stakeholder categories (not individuals) and other sources of requirements
information.}

\section{Environment} \label{sec:environment}

The Environment book describes the domain and external context in which the
system will operate.

\subsection{Glossary} \label{sec:e1}

\wss{Definitions of domain vocabulary, special uses of ordinary words, and
acronyms.}

\subsection{Components} \label{sec:e2}

\wss{External elements that may affect or be affected by the system, including
systems the product must interface with. Distinct from interfaces the system
provides (Section~\ref{sec:s3}) and from technology required to build it
(Section~\ref{sec:p5}).}

\subsection{Constraints} \label{sec:e3}

\wss{Non-negotiable obligations and limits imposed by the environment.}

\subsection{Assumptions} \label{sec:e4}

\wss{Properties assumed to hold in order to simplify construction (not
environment-imposed constraints).}

\subsection{Effects} \label{sec:e5}

\wss{How the system's operation will affect the environment.}

\subsection{Invariants} \label{sec:e6}

\wss{Environment properties the system may assume on entry and must preserve.}

\section{System} \label{sec:system}

The System book refines the Goals book with detailed requirements: structure,
behaviour, interfaces, and acceptance criteria.

\subsection{Components} \label{sec:s1}

\wss{Major software (and, if applicable, hardware) parts.}

\subsection{Functionality} \label{sec:s2}

\wss{Behaviours of the system (functional and non-functional), typically one
subsection per component from Section~\ref{sec:s1}.}

\subsection{Interfaces} \label{sec:s3}

\wss{How functionality is made available to users and other systems (UI, APIs).
Interfaces consumed from outside belong in Section~\ref{sec:e2}.}

\subsection{Detailed Usage Scenarios} \label{sec:s4}

\wss{User stories / scenarios that may reference system components, special
cases, and errors. Complements (does not replace) precise behaviour in
Section~\ref{sec:s2}.}

\subsection{Prioritization} \label{sec:s5}

\wss{Criticality of behaviours, interfaces, and scenarios.}

\subsection{Verification and Acceptance Criteria} \label{sec:s6}

\wss{Conditions under which an implementation is considered satisfactory
(testing and other V\&V).}

\section{Project} \label{sec:project}

The Project book describes the process that produces the system. Where this
overlaps the Development Plan, summarize here and point to that document for
detail.

\subsection{Roles and Personnel} \label{sec:p1}

\wss{Roles needed to carry out the project.}

\subsection{Imposed Technical Choices} \label{sec:p2}

\wss{Technical decisions already imposed (languages, platforms, standards).}

\subsection{Schedule and Milestones} \label{sec:p3}

\wss{Timeline and major milestones.}

\subsection{Tasks and Deliverables} \label{sec:p4}

\wss{Work items and artifacts to produce.}

\subsection{Required Technology Elements} \label{sec:p5}

\wss{Tools and technologies needed to build the system (distinct from
environment components in Section~\ref{sec:e2}).}

\subsection{Risk and Mitigation Analysis} \label{sec:p6}

\wss{Project risks and planned mitigations. Detailed hazard analysis may live
in the Hazard Analysis document.}

\subsection{Requirements Process and Report} \label{sec:p7}

\wss{How requirements were (and will be) elicited, recorded, and maintained.}

\section{References}

\begin{itemize}
  \item Bertrand Meyer. \emph{Handbook of Requirements and Business Analysis}.
  Springer, 2022.
\end{itemize}

\newpage{}

\section{Appendix --- Generative AI Use Disclosure}

\subsection{AI Tools Used}

\wss{Give the name of the AI tool(s) used and the version.}

\subsection{How and Where Used}

\wss{Explain what the tool(s) were used for, and which parts of the document
were affected.}

\subsection{Verification of AI-Generated Content}

\wss{Describe how the team checked AI-generated or AI-assisted material before
including it. The team remains responsible for accuracy and completeness.}

\newpage{}

\section*{Appendix --- Reflection}

\wss{Not required for CAS 741}

\input{../SRS_Reflection.text}

\end{document}
